% ===== §9 Aesthetic Education: From Tradition to the Relational =====
\section{Aesthetic Education: From Tradition to the Relational}
\label{sec:education}

\noindent
The convergence delivered the paper's title as a consequence: because the perceptual frame is generated, plastic, and relationally formed, an education of perception is possible and must be relational. This section begins the development of that education by setting it against the tradition it revises. It is a tradition of the first rank, and the paper approaches it with the respect owed to the thought it is departing from, since the departure is intelligible only against a full statement of what is departed from. The section therefore does three things. It states the tradition of aesthetic education at the level of its theory, from Schiller through Dewey and Cai Yuanpei to Read and Greene, taking care to render each at its strongest and to mark where the tradition already strains toward what this paper will claim. It isolates the shared presupposition that nonetheless unites these otherwise different thinkers, that the unit of aesthetic cultivation is the individual. And it states the difference that follows from the convergence, that generative relational aesthetic education takes the relation for its unit, before showing how the individual and the relational are held together in the cultivation itself. The treatment here is of the tradition's theory; its practical models, the actual methods by which it has trained perception, are taken up later, where they can be set beside the paper's own models of practice (\xref{sec:practicefield}).

\subsection{The tradition at its strongest}
\label{sub:edu-tradition}

Aesthetic education has a founding text, and it set the terms the tradition would keep. Schiller's letters proposed that the human being is divided between a sensuous drive, bound to matter and the life of feeling, and a formal drive, bound to reason and law, and that neither alone yields freedom; the aesthetic, through what he called the play drive, reconciles the two, restoring to the person a wholeness in which sensibility and reason are no longer at war \citep{schiller1967}. It is essential to render Schiller fairly here, because his account is not the private cultivation of taste it is sometimes reduced to. The reconciliation he sought had a political horizon: the aesthetic education of the human being was, for him, the path by which a people might become capable of freedom, the aesthetic state the condition of a genuine political one. Schiller is not a straw individualist, and the paper does not treat him as one. Yet the wholeness his education restores is the wholeness of a person, the reconciliation of the drives within a single soul, and the political freedom it opens onto is reached by way of the aesthetic completion of individuals. The unit is the person, however political the person's completion is meant to become.

The tradition's later high points inherit this unit while enriching the account. Dewey brought the aesthetic down from the museum into the continuity of experience, refusing the separation of art from life and locating the aesthetic in the rhythm of doing and undergoing by which any experience may be brought to fullness; his was a democratic aesthetics, holding that every person is capable of the artist's relation to their own experience \citep{dewey1934}. This is the tradition at its most hospitable to the everyday, and the paper draws on it gratefully. But the artist Dewey would make of every person is each person severally, and the fullness of experience he describes is the fullness of an individual's experience, however much that individual lives among others. Cai Yuanpei, proposing in 1917 that aesthetic education replace religion, appropriated the Kantian and Schillerian account for a China in crisis and gave the tradition one of its most socially ambitious statements, aesthetic education as the cultivation of a citizenry; yet the development he sought was the full development of the person in the three domains of knowledge, will, and feeling, an individually focused vision even at its most national in aim \citep{caiyuanpei1917}. Read, making art the medium of a general education, sought through it the integration of the individual personality \citep{read1943}. And Greene, who of all the tradition comes nearest to what this paper will claim, made aesthetic education the deliberate fostering of informed and imaginative encounters with art, and turned in the end toward community, toward the classroom as a social undertaking in which a teacher must see and hear students in all their plurality \citep{greene1995}. Greene reaches the threshold of the relational and is the tradition's most generous approach to it. But what her aesthetic education releases is the imagination, and the imagination it releases is each student's; the community she envisions is the condition and the medium of that release rather than itself the thing cultivated. The tradition strains toward the relational in Greene and does not cross into it.

\subsection{The shared presupposition}
\label{sub:edu-presupposition}

Across these differences one presupposition holds, and naming it is the point of the survey. In each case the unit of aesthetic cultivation is the individual. What is to be brought toward completion, fullness, integration, or imaginative freedom is a person: Schiller's soul reconciled in play, Dewey's individual brought to the fullness of an experience, Cai's citizen developed in knowledge and will and feeling, Read's integrated personality, Greene's student whose imagination is released. That the completion is often meant to have social or political issue does not alter its unit; the person is completed first, and the polity or the community is reached through or after the completion of persons. This is not a failing of the tradition on its own terms, for the tradition took the cultivation of individuals to be exactly its task. It becomes a limit only in the light of the convergence, which showed that the perceptual frame an aesthetic education would cultivate is not formed in the individual first and shared afterward, but formed in relation from the start. If the frame is relationally constituted, then an education that takes the individual as its unit works downstream of where the frame is actually formed, cultivating the person's sensibility while leaving unaddressed the relation in which that sensibility was and is generated. The tradition's presupposition is not false about individuals; it is incomplete about where perception comes from.

\subsection{The relational difference}
\label{sub:edu-difference}

Generative relational aesthetic education takes the relation for its unit, and this is the whole of its difference from the tradition. What it cultivates is not, in the first instance, an individual's capacity to perceive the everyday aesthetically, but the capacity of a relation to go on generating aesthetic perception between those who stand in it. The distinction is not a matter of emphasis or of adding a social dimension to an individual cultivation. It is a difference in the object of cultivation. Where the tradition asks how a person's sensibility may be brought toward completion, this paper asks how two people may cultivate between them a shared perceptual life that neither carries alone, a capacity resident in the relation rather than in either party, by which their common world is continually generated as perceived. The convergence licenses this shift and indeed requires it: since the frame under which perception is generated is laid down in shared engagement, organised in relation to others' desire, and produced within a social order, the deliberate formation of that frame must work on the relations that form it, and cannot be accomplished by either party working alone on their own sensibility. The everyday is where this is most concretely true. The shared perceptual life of two people who live together, the qualities of light and sound and texture and gesture that their common days are made of, is generated in and by their relation, in the ten thousand small mutual attunements of a shared ordinary life; and to cultivate it is to work on that relation's generative capacity, not to improve two sensibilities in parallel and hope they meet. Relational aesthetic education is the deliberate cultivation of the relation's power to generate the perception of a shared everyday.

\subsection{The individual and the relational in cultivation}
\label{sub:edu-dialectic}

To take the relation as the unit is not to dispense with the individual, and the third face of the paper's thesis, stated at the outset, must now be shown at work in the cultivation itself. The individual and the relational stand in a spiral, and the cultivation runs in both directions of it. An individual's perceptual refinement, the trained discrimination whose neural reality the second spine established, is drawn upon by the relation and is a condition of what the relation can generate: two people cannot generate between them a finely perceived shared world if neither has cultivated the capacity to perceive finely at all, and the sommelier's trained palate, brought into a shared meal, raises what the shared perception of that meal can become. In this direction the individual conditions the relational. But the shared perception generated in the relation returns upon the individual and reshapes their capacity: what two people come to perceive together, each learns to perceive, so that the individual enters the next occasion already altered by what the relation generated, perceiving now what before the relation they could not. In this direction the relational conditions the individual. Neither pole is prior, and the cultivation is the turning of this spiral, in which refined individual perception feeds the relational generation and the relational generation refines individual perception, and the return does not arrive back at its origin, since each turn leaves both the individuals and the relation altered from where they began. This is why relational aesthetic education does not swallow the individual it declines to take as its unit. It cultivates the relation, and in cultivating the relation it cultivates the individuals, because the relation and the individuals are joined in a spiral that the cultivation sets and keeps turning. The tradition cultivated individuals and hoped the relation would follow; this paper cultivates the relation and shows that the individuals follow, because the perceptual frame was relational before it was ever individual, and returns to the individual by way of the relation that formed it.
